\chapter{Charakterisierungen und Randfälle}
\label{b06guide:ch:criteria}

\section{Punkte und Fasern}

\begin{theorem}[Punkt- und Faserkriterien]
\label{b06guide:thm:pointcriteria}
Für eine Funktion $f\colon A\to B$ sind äquivalent:
\begin{enumerate}
\item $f$ ist injektiv.
\item Für alle $x,y\in A$ gilt
      $f(x)=f(y)\Longleftrightarrow x=y$.
\item Für alle $x,y\in A$ gilt
      $x\ne y\Rightarrow f(x)\ne f(y)$.
\item Jede Faser $\operatorname{Fib}_f(b)$, $b\in B$,
      enthält höchstens ein Element.
\item Für jedes $b\in f[A]$ gibt es genau ein $x\in A$ mit $f(x)=b$.
\end{enumerate}
Fasern über Elementen von $B\setminus f[A]$ dürfen leer sein.
\end{theorem}

\begin{theorem}[Prüfung auf höchstens zweielementigen Teilmengen]
\label{b06guide:thm:twopoints}
Eine Funktion $f\colon A\to B$ ist genau dann injektiv,
wenn für alle $x,y\in A$ die Einschränkung auf $\{x,y\}$
injektiv ist. Gleichwertig ist die Forderung
\[
 x,y\in A,\quad x\ne y
 \quad\Longrightarrow\quad f[\{x\}]\cap f[\{y\}]=\varnothing.
\]
\end{theorem}

\section{Charakterisierung durch Teilmengen}

\begin{definition}[Bildfunktion auf der Potenzmenge]
Für $f\colon A\to B$ sei
\[
 f_*\colon\mathcal P(A)\to\mathcal P(B),\qquad U\mapsto f[U].
\]
\end{definition}

\begin{theorem}[Rückgewinnung von Teilmengen]
\label{b06guide:thm:recoverycriteria}
Für $f\colon A\to B$ sind äquivalent:
\begin{enumerate}
\item $f$ ist injektiv.
\item Für jede Teilmenge $U\subseteq A$ gilt $f^{-1}[f[U]]=U$.
\item Die Bildfunktion $f_*$ ist injektiv.
\item Für alle $U,W\subseteq A$ gilt
      $f[U]\subseteq f[W]\Rightarrow U\subseteq W$.
\end{enumerate}
In Bedingung~2 genügt bereits die Prüfung aller Einermengen
$U=\{x\}$ mit $x\in A$.
\end{theorem}

\begin{theorem}[Charakterisierung durch Mengenoperationen]
\label{b06guide:thm:operationcriteria}
Für $f\colon A\to B$ ist jede der folgenden Bedingungen
für sich gleichwertig mit der Injektivität von $f$:
\begin{enumerate}
\item Für alle $U,W\subseteq A$ gilt
      $f[U\cap W]=f[U]\cap f[W]$.
\item Für alle disjunkten $U,W\subseteq A$ sind
      auch $f[U]$ und $f[W]$ disjunkt.
\item Für alle $U,W\subseteq A$ gilt
      $f[U\setminus W]=f[U]\setminus f[W]$.
\item Für alle $U\subseteq A$ gilt
      $f[A\setminus U]=f[A]\setminus f[U]$.
\item Für alle $U,W\subseteq A$ gilt
      $f[U\mathbin{\triangle}W]=f[U]\mathbin{\triangle}f[W]$.
\end{enumerate}
Die Forderungen gelten jeweils für \emph{alle} genannten
Teilmengen. Eine einzelne zufällig erfüllte Gleichheit genügt nicht.
\end{theorem}

\section{Leere und kleine Definitionsbereiche}

\begin{theorem}[Leere Menge und Einermengen]
\label{b06guide:thm:empty}
Für jede Menge $B$ gilt
\[
 \operatorname{Inj}(\varnothing,B)=\{\varnothing\}.
\]
Jede Funktion mit höchstens einem Element im Definitionsbereich
ist injektiv. Insbesondere ist für $b\in B$ die Funktion
$\{(a,b)\}\colon\{a\}\to B$ injektiv.
\end{theorem}

\begin{theorem}[Konstante Funktionen und kleine Zielmengen]
\label{b06guide:thm:constant}
Seien $b\in B$ und $c\colon A\to B$ mit $c(x)=b$ für alle $x\in A$.
Dann ist $c$ genau dann injektiv, wenn $A$ höchstens ein Element
enthält. Insbesondere gilt
\[
 \begin{aligned}
 \operatorname{Inj}(A,\{b\})\ne\varnothing
   &\Longleftrightarrow A\text{ enthält höchstens ein Element},\\
 \operatorname{Inj}(A,\varnothing)\ne\varnothing
   &\Longleftrightarrow A=\varnothing.
 \end{aligned}
\]
\end{theorem}